\section{导子的泛性质:非交换情形}


记号约定:
\begin{itemize}
	\item 本节中,所提到的代数都是含幺结合代数,代数同态都是含幺代数同态.
	\item $A$是$K$-代数,$A\otimes_{K}A$配备典范的$\left(A,A\right)$-双模结构.
	\item 定义了$A$上乘法的$K$-线性映射$m:A\otimes_{K}A\to A$是$\left(A,A\right)$-双模同态,$I=\ker\left(m\right)$.
	\item $\delta_{A}:A\to I,x\mapsto x\otimes 1_{A}-1_{A}\otimes x$.
\end{itemize}

\begin{lemma}{典范映射$\delta_{A}$是导子}{}
	$\delta_{A}$是$K$-导子,并且$I$作为左$A$-模由$\im\left(\delta_{A}\right)$生成.
\end{lemma}

\begin{proposition}{典范导子的泛性质}{}
	对每个$\left(A,A\right)$-双模$E$和每个$K$-导子$d:A\to E$,存在唯一的$\left(A,A\right)$-双模同态$f:I\to E$使下图交换.
	\begin{center}
		\begin{tikzcd}
			A && E \\
			& I
			\arrow["f", from=1-1, to=1-3]
			\arrow["{\delta_{A}}"', from=1-1, to=2-2]
			\arrow["{\exists!f}"', dashed, from=2-2, to=1-3]
		\end{tikzcd}
	\end{center}
\end{proposition}

\begin{remark}{}{}
	上述命题给出典范$K$-模同构$\Hom_{\left(A,A\right)}\left(I,E\right)\to \Der_{K}\left(A,E\right),f\mapsto f\circ\delta_{A}$.
\end{remark}